\section{Differentiation and numerical contracts}
For fixed relation indices, weighted sum has a simple adjoint. With output
cotangent denoted by a bar, source and edge derivatives are
$$
 \bar{x}_j=\sum_{e:s(e)=j}w_e\bar{y}_{d(e)},\qquad
 \bar{w}_e=x_{s(e)}\bar{y}_{d(e)}.
$$
Repeated source indices require accumulation, not assignment. Empty destination
rows contribute no source/edge derivative. The example in Section~\ref{sec:programming-model} provides
an independently checkable forward and backward value contract.

\subsection{Native and Torch paths}
Native reverse-mode transforms construct derivative expressions from supported
Tensor and relation operations. Torch-facing calls preserve ordinary Torch
outputs and autograd integration, but their backward implementation is path-specific.
Some compiled CSR forwards replay fixed-topology Torch semantics during backward.
That route is differentiable without being a compiled TTIR backward kernel.
Explicitly wrapping a Torch tensor into a native tensor shares storage but does
not merge the two autograd histories.

Supported traced edge neural networks have a CUDA tile forward and a symbolic-VJP
recompute kernel. The tests compare output and gradients of source features,
positions and every module parameter to an eager reference. CPU tests additionally exercise explicit edge
attributes with double-precision finite differences. These checks establish the
tested first-order envelope, not arbitrary Python capture or universal
higher-order differentiation. ReLU-like nondifferentiabilities need appropriate
test points, rather than unqualified finite-difference equality at a kink.

\subsection{Online softmax and approximation}
For messages consisting of a score and value, the streaming state tracks the
maximum score, normalized mass and weighted value accumulator. Two nonempty
states combine by rescaling into a common maximum:
$$
 m=\max(m_a,m_b),\qquad
 \ell=e^{m_a-m}\ell_a+e^{m_b-m}\ell_b,\qquad
 o=e^{m_a-m}o_a+e^{m_b-m}o_b.
$$
Finalization divides the accumulator by the mass for nonempty rows; empty-state
handling must avoid undefined subtraction of infinities and follow the
implementation's reducer convention. This numerical structure is related to
IO-aware attention algorithms~\cite{flashattention}; it does not imply that every
Tiga path achieves their performance or full feature coverage.

An optional positive pruning threshold enables an approximate dense-attention
forward path. Scores and key reads are still needed for the pruning decision;
eligible tiles skip value loading and accumulation. A large score gap can make
a tile's contribution small, but a block-level criterion is not a per-query
error bound. Threshold choice affects accuracy and work saved, and may not improve
latency. The exact default is a disabled threshold, not a threshold of zero.
Pruned backward is not claimed. Exact and approximate workloads must not be mixed
in a single speedup comparison without reporting numerical error.
